% !TeX root = ../Book1.tex
\begin{figure}[htbp]
	\centering
	\begin{tikzpicture}[font=\small]
		\begin{scope}[xshift=-2.8cm,x=2cm,y=1.3cm]
			\draw[->,GeometryGray] (-1.25,0)--(1.3,0) node[right] {\(x\)};
			\draw[->,GeometryGray] (0,-1.2)--(0,1.25) node[above] {\(y\)};
			\draw[GeometryGray,thick,domain=-180:180,samples=181,variable=\t] plot ({sin(\t)},{sin(2*\t)});
			\draw[GeometryBlue,very thick,domain=-27:27,samples=40,variable=\t] plot ({sin(\t)},{sin(2*\t)});
			\draw[orange!75!black,very thick,domain=145:179.8,samples=35,variable=\t] plot ({sin(\t)},{sin(2*\t)});
			\foreach \t in {150,165,174,178}{\fill[orange!75!black] ({sin(\t)},{sin(2*\t)}) circle (1.25pt);}
			\fill[GeometryBlue] (0,0) circle (1.6pt);
			\node[above left] at (0,0) {\(\gamma(0)\)};
			\node[GeometryBlue,above right] at (0.5,0.78) {\(t\approx0\)};
			\node[orange!75!black,right] at (0.63,-0.72) {\(t\uparrow\pi\)};
		\end{scope}
		\begin{scope}[xshift=3.7cm,yshift=0.2cm]
			\draw[GeometryGray] (-1.7,0)--(1.7,0);
			\draw[fill=white] (-1.7,0) circle (2pt);
			\draw[fill=white] (1.7,0) circle (2pt);
			\fill[GeometryBlue] (0,0) circle (2pt);
			\node[below] at (-1.7,0) {\(-\pi\)};
			\node[below] at (0,0) {\(0\)};
			\node[below] at (1.7,0) {\(\pi\)};
			\foreach \x in {1,1.35,1.52,1.61}{\fill[orange!75!black] (\x,0) circle (1.5pt);}
			\draw[->,orange!75!black] (0.95,0.35)--(1.65,0.35);
			\node[above] at (0,0.75) {参数区间 \((-\pi,\pi)\)};
			\node[align=center] at (0,-1) {\(\gamma(t_j)\to\gamma(0)\)\\但 \(t_j\not\to0\)};
		\end{scope}
	\end{tikzpicture}
	\caption[单射浸入的像拓扑]{曲线 \(\gamma(t)=(\sin t,\sin2t)\)，定义域为 \((-\pi,\pi)\)。蓝色弧是零附近的参数片；橙色点来自趋向 \(\pi\) 的参数，也在像中趋向原点。逆参数映射因而不连续，尽管每个参数点的导数非零且参数映射单射。}
	\label{g1:fig:injective-immersion-topology}
\end{figure}
